\section{Optimisation of Telescopes Pointing and Movement in Radio Astronomy}
The precision with which a radio telescope maintains tracking and pointing on celestial sources is the very essence of the integrity and scientific value of data it collects. Pointing or tracking inaccuracies, even in the range of arcseconds, can have significant impacts on detecting faint cosmic signals or lower the quality of interferometric imaging. As a result, the optimisation of the motion and position of telescopes has emerged as an essential research direction in observational astronomy. The area encompasses diverse technologies and strategies such as developing high-performance servo control systems, real-time error correction architectures, adaptive tracking algorithms, and end-to-end software architectures.

The recent advances in technology have also enabled telescopes to be increasingly autonomous, faster-acting, and more accurate in dynamic conditions. The conjunction of predictive control methods, stable error modeling packages, and machine-learning-based back-up systems has now enabled in-time compensation of environmental perturbations, mechanical non-linearities, and structure flexure in real time. Such technologies play particularly important roles in modern-day radio observatories, where large-duration observation, extensive-sky surveys, and high-accuracy timing are routinely demanded.

This subsection covers fundamental principles and mechanisms employed to optimise telescope motion and pointing. The chapter begins with a general overview of motion control approaches, followed by detailed discussion of error compensation models, trajectory generation techniques, servo hardware integration, robotic automation, artificial intelligence-based tracking systems, and integrated control architectures. These components comprise the technological cornerstone of precision radio astronomy and co-evolve with increasing scientific demands on contemporary telescopes.

    \subsection{Principles of Telescope Motion and Control}
    The precise motion control of a radio telescope is central to the success of modern observing astronomy, especially for faint radio sources that need high fidelity pointing for extended periods. Two fundamental functions of telescope control, pointing and tracking, often are performed within integrated real-time servo systems. Slewing to a specified coordinate in equatorial or alt-azimuth forms is the act of pointing. Once the initial orientation is established, tracking is employed to remain in synchronisation with a celestial object as it crosses the sky due to Earth's rotation because. Such operations require the motors, encoders, and sensors of the telescope to work in coordination under the guidance of a master control program.

    Traditionally, PID (Proportional-Integral-Derivative) controllers have been employed for telescope actuation because they are simple and efficient in eliminating steady-state error. PID controllers operate by constantly modifying actuator output to reduce the difference between desired and actual position. PID systems, however, may have difficulty with dynamic disturbances, particularly in large radio telescopes where mechanical backlash, thermal expansion, and wind can cause sudden changes in error. Due to this, researchers have moved towards more advanced control methods, including LQG (Linear Quadratic Gaussian) and H-infinity ($H_\infty$) control frameworks. These methods provide optimal or robust solutions for disturbance elimination and system modeling uncertainties.

    LQG controllers blend state estimation with Kalman filters and optimal control law calculation to optimise a cost function that weighs tracking accuracy against control effort. In contrast, $H_\infty$ controllers pursue robustness with a guaranteed bound on system response in worst-case disturbance. These techniques are important for large-diameter, high-gain antennas operating under environmental stressors, as \cite{Gawronski2007} identified in his groundbreaking work on antenna control. Their use provides enhanced control stability, precision, and faster target acquisition even in changing operating conditions \cite{Gawronski2007}.

    \subsection{Pointing Error Correction and Kinematic Modeling}
    Correct pointing of telescopes is not merely a function of instructing the telescope to point at some celestial coordinates. In practice, there are numerous types of errors that build up from structural, mechanical, and environmental origins that have to be modeled and adjusted constantly in order to preserve observation accuracy. Some of these include axis misalignment errors, encoder drift, flexure caused by gravitational loading, thermal expansion, and atmospheric refraction. Therefore, pointing systems should be designed with strong correction models that incorporate both real-time feedback and historical calibration data to achieve accurate tracking and initial acquisition.

    Pointing errors are generally of two types: systematic errors, which are repeatable and predictable, and random or dynamic errors, which vary in time due to environmental effects. Systematic errors are compensated for with the help of telescopes' kinematic calibration models—mathematical representations of the telescope geometry and mechanical imperfections. One of these programs, which has seen widespread use in the astronomical community, is TPOINT, a modeling program that utilises regression analysis to detect and correct for more than 20 frequent sources of error, such as non-perpendicular axes, mechanical flexure, and encoder offsets. TPOINT implements such corrections as fitted parameters, which are merely subtracted from raw position commands, thereby improving the intrinsic accuracy of the telescope.

    While traditional pointing models like TPOINT are sufficient in most cases, they may not fare well in rapidly changing conditions or high-velocity tracking situations. For this, recent advancements have adopted machine learning and AI-based correction systems that learn in real-time. \cite{He2023, Waltenberger2024} suggested a hybrid model that blends TPOINT's statistical modeling with reinforcement learning methods to develop a more robust correction model that can adapt to time-varying errors. Similarly, He et al. (2023) developed an optimised parameter model that used both real-time encoder feedback and pre-trained kinematic data to significantly reduce dynamic pointing error. These systems represent a new generation of adaptive control in radio astronomy, blending physics-based and data-driven methodologies \cite{He2023, Waltenberger2024}.

    \subsection{Trajectory Planning and Control Algorithms}
    Trajectory planning for telescope motion is an essential part in enabling telescope slews to be fast, smooth, and mechanically safe. Trajectory planning here means the computation of the path that the telescope must follow to move between two pointing positions. The process must consider limits on maximum velocity, acceleration, and mechanical limits of the telescope structure. If improperly planned trajectories, fast or jerky motion can excite mechanical resonances, wear and tear upon actuators and gears, and vibrates to compromise image quality or interfere with radio signal reception. Standardly, telescope motion profiles simply relied on simple velocity ramps or trapezoidal profiles, assuming strictly constant acceleration and deceleration phases. While good enough in many applications, however, such approaches tend to induce unwanted mechanical stresses.

    To surpass these limitations, even more sophisticated methods have been developed. One of these methods is to use jerk-limited polynomial trajectories, where not only is the acceleration limited, but so too is the time derivative of the acceleration—the so-called "jerk". This results in even smoother operations that significantly reduce the excitation of structural modes, thereby ensuring alignment and structural integrity during slewing. \cite{Savarese2020} described that the approach allows rapid yet stable repositioning and, as a result, it is ideal for surveillance programs with high target change rates \cite{Savarese2020}.

    Besides trajectory smoothing, modern control techniques also include Model Predictive Control (MPC), which dynamically recalculates optimum motion path in real time according to the current state of the telescope and expected future behavior. This enables adaptive compensation for wind loads, mechanical backlash, or drive irregularities. For example, \cite{Petrillo2024} applied MPC to the Telescopio Nazionale Galileo and reported improved performance relative to conventional PID controllers under real observation conditions. MPC offered improved damping, reduced overshoot, and closer target coordinate convergence even in the presence of changing environmental disturbances \cite{Petrillo2024}. These findings show the increasingly prominent role predictive, optimisation-based algorithms will play in future-generation telescope motion planning.

    \subsection{Servo Control and Hardware Integration Enhancements}
    Servo control systems form the heart of precision telescope motion. They constitute the interface between high-level trajectory commands and physical actuation of the telescope's axes. For large-scale radio telescopes, the complexity of this control is greatly increased due to structural flexibility, gear backlash, wind-induced vibrations, and thermal distortions. As a result, the servo system not only has to follow position commands but also must actively cancel dynamic disturbances that can degrade pointing stability. Integration of control software with mechanical and electronic hardware is thus a key design issue, one which demands intimate knowledge of system dynamics and feedback characteristics.

    Traditional PID-type servo systems are not suitable for large, flexible structures because they are not predictive and have limitations in maintaining robustness in the presence of varying loads or system nonlinearities. To overcome such shortcomings, researchers have formulated state-space controllers, and more precisely, the Linear Quadratic Gaussian (LQG) technique, which combines optimum control with real-time state estimation. At the Telescopio Nazionale Galileo (TNG), \cite{Schipani2020} commissioned an LQG-based servo system that had good mechanical backlash suppression and stability against structural oscillations. The system employed real-time sensor fusion and internal model prediction for suppressing vibrational feedback, with a performance better than traditional methods \cite{Schipani2020}.

    Additional progress has since been achieved through the implementation of two-degree-of-freedom (2-DOF) controllers, more specifically based on H-infinity ($H_\infty$) robust control. Such controllers split the reference tracking performance from the disturbance rejection function so that there can be more specialised tuning for each response. \cite{Yang2023} applied a 2-DOF ($H_\infty$) controller in a big radio telescope servo system and realised great improvement in tracking accuracy and resistance to wind interference and mechanical noise. The method allows for better control in the worst-case scenario and improves telescope availability by maintaining stable performance under harsh environmental conditions \cite{Yang2023}. The marriage of such powerful control techniques with hard real-time back-ends to hardware is an important advance in the design of telescope control systems.

    \subsection{Robotic and Automation Control}
    The enhanced complexity and operating demands of today's observatories have pushed the trend towards automatic and robotic telescope systems, particularly for small- to medium-scale research centers and global networks of telescopes. These systems will perform all of the observational activities autonomously with no intervention on the part of humans, i.e., telescope slewing, tracking, focusing, environmental monitoring, data acquisition, and shut-down processes. The main purpose of robotic control is to maximise observational uptime with a minimum requirement for on-site personnel, especially in remote or high-altitude locations where human interaction and maintenance cost are minimal.

    At the core of robotic telescope systems is an architecture of closed-loop control supporting real-time feedback and self-healing. Its feedback loop permanently cross-checks the telescope's current position and its commanded course, and adaptively adjusts the motor commands so as to continuously maintain it perfectly aligned. Its environment sensors including anemometers, temperature probes, and cloud detectors feed in additional information into the control system so that the system can autonomously react to environmental changes, like halting activities during bad weather. This integration of environmental information adds a further level of fault tolerance and safety necessary for unsupervised monitoring.

    A good example of such implementation is the 50 cm robotic telescope system which was refurbished by \cite{Stanzin2022}. They refurbished the telescope by fitting it with advanced automation routines and closed-loop pointing corrections to enable it to run in a completely unattended mode. The automation routines within the system made it possible for easy reconfiguration according to newer observational priorities as well as environmental inputs. Besides, the internal processes of the controller managed self-calibration, target acquisition, and health diagnostics to offer uniform performance during long observing sessions \cite{Stanzin2022}. The above innovations reflect how robot control has transformed telescope operations from operator-controlled systems to intelligent platforms for high-efficiency autonomous scientific productivity.

    \subsection{Predictive Models Using AI to Optimise Tracking}
    As radio telescopes get more accurate and observation cadence increases, traditional reactive control techniques are being augmented—or even substituted—by predictive models based on artificial intelligence (AI). They constitute a new paradigm for controlling telescopes: instead of responding to tracking errors once they have occurred, AI-based systems forecast and prevent them by learning from past performance data and anticipating future system behavior. This kind of method is highly practical for telescopes involved in high-speed tracking, dynamic scheduling of observations, or multi-object observing campaigns where strict alignment and fast repositioning are absolutely imperative.

    At the center of such systems are algorithms that perform online learning, i.e., wherein the model constantly updates its internal parameters automatically in real-time using new data input without requiring offline retraining. This makes the control system able to learn to respond to changes in behavior of the telescope due to mechanical wear, thermal variations, seasonal sagging of structures, or unexpected environmental disturbances. Reinforcement learning—a type of AI in which the model is trained to act optimally by trial and error in an environment, real or simulated—is being used more and more in telescope motion control.

    \cite{Waltenberger2024} depicted reinforcement learning of telescopic high-speed tracking improvement. Their method was trained to minimise pointing jitter through pre-correction of motion with a hybrid architecture of deterministic physics-based models augmented by data-driven prediction modules. This allowed adaptive correction of non-linear dynamics hard to analyse manually \cite{Waltenberger2024}. Unlike fixed-parameter controllers, AI-enhanced systems are able to learn over time and improve their performance as they encounter more operating conditions. They can identify patterns and anomalies, enabling predictive maintenance by detecting the telltale early warning signs of actuator degradation or alignment drift. Those capabilities are paving the way toward a future generation of smart telescope control architectures that combine real-time response with strategic foresight.

    \subsection{Unified Control System Architecture}
    In more complex observatories, notably those with multiple instruments, reused infrastructure, or distributed antenna systems, the need for a unified control system architecture has reached top priority. A unified system integrates all subsystems—azimuth and elevation motion control, dome rotation, environmental monitoring, power distribution, data acquisition, and time synchronisation—into a single homogeneous software and hardware platform. This offers consistency of operation, minimises latency between modules, and allows centralised error processing, diagnostics, and system optimisation.

    One of the key applications for such an architecture in radio astronomy is synchronisation, both pointing telescope and timing precision. This is particularly applicable to interferometric arrays, where signals from several telescopes must be coherently combined. Such systems need nanosecond-level synchronisation, typically achieved using GPS-disciplined oscillators, atomic clocks, and precise network time protocols. The control software must not only regulate synchronisation at the hardware level but also synchronise observation scheduling, real-time data correlation, and system health monitoring between geographically dispersed components.

    Integrated systems also facilitate the updating of legacy hardware. Where decommissioned satellite communications antennas, for instance, are reused as working radio telescopes, control architecture must bridge digital command protocols and analog or outdated hardware. \cite{Ramadhani2022} accomplished this in their retrofitting of a 7.3-meter telecommunication antenna as a scientifically functional radio telescope. Their approach was the addition of environmental feedback systems, precise azimuth-elevation control, and a networked timing module to an integrated software environment. This provided precise tracking, improved reliability, and for the instrument to be remotely operated in entirety \cite{Ramadhani2022}.

    On top of that, integrated architecture is scalable. As one expands or adds more instruments to observatories, modular design enables one to add or substitute parts without requiring a redesign of the entire system. By combining mechanical actuation, sensor integration, data management, and computational control, an integrated control system is the foundation of modern radio astronomy infrastructure that facilitates continuous innovation and seamless international collaboration.